\par\noindent\begin{minipage}{\linewidth}
% !TeX root = ../教材题库_全部条目_转写工作稿.tex
% Source: XJ-B1-C01-HW-P023-Y0993
% 第1章 / 1.2 常用逻辑用语 / 教材 p.23 / PDF 第30页
% 已对照原页转写题干及公式；待教师复核；未制作答案
% 原题及相邻原解截图（核对用）：../题目原图/XJ-B1-C01-HW-P023-Y0993-01.webp
% 本文件只在转写完成后作为题干使用；OCR 条目尚不是可靠题干。
\begin{bankitem}{习题1.2 第2题}{第1章 / 1.2 常用逻辑用语 / 教材 p.23 / PDF 第30页}{XJ-B1-C01-HW-P023-Y0993}{已对照原页转写题干及公式；待教师复核；未制作答案}
写出下列命题的否定，并判断其真假。
\subq{1}$p$：边长分别为 5，12，13 的三角形是直角三角形；
\subq{2}$q$：面积相等的三角形都是全等三角形；
\subq{3}$s$：二次函数 $y=x^2-2x+3$ 的图象与 $x$ 轴没有交点。
\end{bankitem}
\end{minipage}\par\medskip
